\documentclass[11pt,a4paper]{article}
\usepackage[utf8]{inputenc}
\usepackage[portuguese]{babel}
\usepackage{geometry}
\usepackage{graphicx}
\usepackage{titlesec}
\usepackage{parskip}
\usepackage{listings}
\usepackage{xcolor}
\usepackage{float}
\usepackage{hyperref}
\usepackage{enumitem}
\usepackage{amsmath}

\geometry{margin=2.5cm}

% Configuração de código
\lstset{
    basicstyle=\small\ttfamily,
    breaklines=true,
    frame=single,
    numbers=left,
    numberstyle=\tiny,
    commentstyle=\color{gray},
    keywordstyle=\color{blue},
    stringstyle=\color{red},
    language=Python
}

% Formatação de títulos
\titleformat{\section}{\Large\bfseries}{\thesection}{1em}{}
\titleformat{\subsection}{\large\bfseries}{\thesubsection}{1em}{}

\begin{document}

% PÁGINA DE CAPA
\begin{titlepage}
    \centering
    
    \includegraphics[width=0.3\textwidth]{uminhoLogo.png}\\[0.5cm]
    
    {\large \textbf{Universidade do Minho}}\\[0.3cm]
    {\large Escola de Engenharia}\\[0.2cm]
    {\large Departamento de Informática}\\[3cm]
    
    {\Huge \textbf{TaxiGreen:\\[0.3cm]Gestão Inteligente de Frota}}\\[1cm]
    
    {\Large Inteligência Artificial}\\[0.3cm]
    {\large Trabalho Prático de Grupo}\\[0.3cm]
    {\large Ano Letivo 2024/2025}\\[3cm]
    
    \begin{flushright}
        \textbf{Grupo 18}\\[0.5cm]
        Francisco Soares - A106882\\
        Carlos Silva - A106792\\
        Nuno Soares - A107366\\
        Carolina Gonçalves - A106897\\
    \end{flushright}
    
    \vfill
    
    {\large Janeiro de 2026}
    
\end{titlepage}

% PÁGINA DE AVALIAÇÃO PELOS PARES
\section*{Avaliação pelos Pares}

Após análise coletiva do contributo de cada elemento do grupo, apresentamos os deltas correspondentes à componente de avaliação pelos pares. Esta avaliação reflete o esforço e dedicação individual de cada membro ao longo do desenvolvimento do trabalho.

\begin{center}
\begin{tabular}{|l|l|c|}
\hline
\textbf{Número} & \textbf{Nome} & \textbf{DELTA} \\
\hline
A106882 & Francisco Soares & +0.15 \\
A106792 & Carlos Silva & +0.15 \\
A107366 & Nuno Soares & +0.15 \\
A106897 & Carolina Gonçalves & -0.45 \\
\hline
\end{tabular}
\end{center}

\textbf{Nota:} O somatório dos DELTAs é igual a 0.00.

\newpage

% ÍNDICE
\tableofcontents
\newpage

% CONTEÚDO DO RELATÓRIO
\section{Introdução}

Neste trabalho, o nosso grupo propôs-se a desenvolver um sistema de gestão inteligente para a frota da TaxiGreen, uma empresa de táxis urbanos que opera com veículos elétricos e a combustão. Reconhecemos que a gestão eficiente de uma frota heterogénea representa um desafio complexo, especialmente quando consideramos a necessidade de equilibrar custos operacionais, tempos de resposta e impacto ambiental.

A TaxiGreen enfrenta desafios crescentes na sua operação diária. Por um lado, os veículos elétricos oferecem custos operacionais reduzidos e zero emissões locais, mas apresentam autonomia limitada e tempos de recarga significativos. Por outro lado, os veículos a combustão possuem maior flexibilidade operacional, embora com custos e impacto ambiental superiores. O nosso sistema pretende encontrar o equilíbrio ideal entre estas características distintas.

O objetivo principal deste trabalho consiste em desenvolver algoritmos de procura que permitam otimizar a gestão da frota, garantindo eficiência operacional, redução de custos energéticos e cumprimento de critérios ambientais. Ao longo deste relatório, apresentamos a formulação do problema como um problema de procura, descrevemos os algoritmos implementados, analisamos as estratégias de alocação desenvolvidas e discutimos os resultados obtidos nas simulações realizadas.

\section{Descrição do Problema}

A TaxiGreen opera uma frota mista de veículos elétricos e a combustão, cada tipo com características distintas. Os veículos elétricos têm autonomia de 150-300 km, custo de 0.10-0.15 EUR/km, recarga de 30-60 min e zero emissões. Os veículos a combustão possuem autonomia de 500-700 km, custo de 0.20-0.30 EUR/km, reabastecimento de 5-10 min e emissões de 120-180 g CO$_2$/km.

O problema é caracterizado por pedidos dinâmicos ao longo do dia, condições de trânsito variáveis e eventos imprevistos (estradas cortadas, falhas em estações). Os principais desafios incluem responder rapidamente aos pedidos, gerir veículos heterogéneos, otimizar custos operacionais e reduzir emissões de CO$_2$.

Os pedidos de transporte especificam localização de origem e destino, número de passageiros, horário pretendido, nível de prioridade (normal, urgente) e preferência ambiental do cliente. Este conjunto de requisitos torna o problema particularmente complexo, exigindo uma abordagem que considere múltiplos critérios de otimização simultaneamente.

\section{Formulação do Problema}

Representamos o estado através de um tuplo $S = (V, P, C, T)$, onde $V$ é o conjunto de veículos com propriedades como identificador, tipo, localização, autonomia atual/máxima, capacidade e custo por km. $P$ representa os pedidos pendentes com origem, destino, passageiros, horário e prioridade. $C$ descreve o estado da cidade incluindo grafo de localizações, estações e fatores de trânsito. $T$ indica o tempo atual da simulação.

No estado inicial, todos os veículos estão disponíveis com autonomia máxima, posicionados em localizações estratégicas. A cidade opera em condições normais com trânsito a fator 1.0, todas as estradas abertas e estações funcionais.

O objetivo é satisfazer continuamente os pedidos maximizando a taxa de atendimento, minimizando tempo de resposta, custos operacionais, emissões de CO$_2$ e quilómetros vazios (percorridos sem passageiros).

Os operadores definidos são:
\begin{itemize}
    \item \texttt{ATRIBUIR\_VEICULO(v, p)}: associa veículo a pedido (pré-condições: disponível, capacidade, autonomia)
    \item \texttt{DESLOCAR\_VEICULO(v, origem, destino)}: move veículo (consome autonomia, tempo)
    \item \texttt{RECARREGAR(v, e)}: repõe autonomia em estação compatível
    \item \texttt{REJEITAR\_PEDIDO(p)}: para casos inviáveis
    \item \texttt{AGUARDAR(v, duracao)}: veículo permanece em espera
\end{itemize}

A função de custo multi-objetivo pondera diferentes critérios: $C_{total} = \alpha_1 \cdot C_{operacional} + \alpha_2 \cdot C_{tempo} + \alpha_3 \cdot C_{ambiental} + \alpha_4 \cdot C_{penalizacoes}$, onde os pesos $\alpha$ são configuráveis conforme as prioridades operacionais.

\section{Representação da Cidade}

A cidade é representada como grafo direcionado $G = (N, E)$ onde os nós correspondem a localizações relevantes e as arestas representam ligações rodoviárias. Os tipos de nós incluem:
\begin{itemize}
    \item Zonas de recolha de passageiros distribuídas pela cidade
    \item Estações de recarga elétrica para veículos elétricos
    \item Postos de abastecimento para veículos a combustão
\end{itemize}

Cada aresta possui atributos de distância em quilómetros, tempo base de viagem calculado a partir da velocidade média, fator de trânsito atual (multiplicador do tempo base) e estado da via (aberta ou cortada).

Esta representação permite modelar condições dinâmicas que afetam a operação:
\begin{itemize}
    \item Congestionamento de trânsito: fatores entre 1.0 (livre) e 3.0 (severo)
    \item Estradas cortadas: arestas intransitáveis forçando rotas alternativas
    \item Disponibilidade de estações: simulando avarias ou manutenção
\end{itemize}

\section{Algoritmos de Procura}

\subsection{Algoritmos Não-Informados}

Implementámos três algoritmos de procura não-informada que exploram o espaço de estados sem utilizar informação específica do domínio:

\textbf{BFS (Breadth-First Search)} explora todos os nós a uma determinada profundidade antes de avançar, utilizando uma fila FIFO. É completo e ótimo para custos uniformes, com complexidade temporal e espacial $O(b^d)$ onde $b$ é o fator de ramificação e $d$ a profundidade da solução.

\textbf{DFS (Depth-First Search)} explora cada ramo até à máxima profundidade antes de retroceder, utilizando uma pilha LIFO. Tem menor consumo de memória com complexidade espacial $O(bm)$ onde $m$ é a profundidade máxima, mas não garante soluções ótimas.

\textbf{UCS (Uniform Cost Search)} expande sempre o nó com menor custo acumulado $g(n)$, utilizando uma fila de prioridade. Garante otimalidade independentemente da variação de custos entre ações, sendo adequado quando os custos variam conforme veículo e trânsito.

\subsection{Algoritmos Informados}

Desenvolvemos algoritmos que utilizam heurísticas para guiar a procura de forma mais eficiente:

\textbf{A*} combina o custo do caminho $g(n)$ com estimativa heurística $h(n)$, expandindo por ordem de $f(n) = g(n) + h(n)$. É completo e ótimo quando a heurística é admissível (nunca sobrestima o custo real). A eficiência depende diretamente da qualidade da heurística.

\textbf{Greedy Best-First Search} expande o nó com menor valor heurístico $h(n)$, ignorando o custo acumulado. É tipicamente mais rápido pois foca-se no objetivo, mas não garante soluções ótimas e pode ficar presa em mínimos locais.

\subsection{Comparação de Complexidade}

\begin{center}
\begin{tabular}{|l|c|c|c|c|}
\hline
\textbf{Algoritmo} & \textbf{Completo} & \textbf{Ótimo} & \textbf{Tempo} & \textbf{Espaço} \\
\hline
BFS & Sim & Sim* & $O(b^d)$ & $O(b^d)$ \\
DFS & Não & Não & $O(b^m)$ & $O(bm)$ \\
UCS & Sim & Sim & $O(b^{C^*/\epsilon})$ & $O(b^{C^*/\epsilon})$ \\
A* & Sim & Sim** & $O(b^d)$ & $O(b^d)$ \\
Greedy & Não & Não & $O(b^m)$ & $O(b^m)$ \\
\hline
\end{tabular}
\end{center}
\textit{* Para custos uniformes. ** Para heurísticas admissíveis.}

\subsection{Heurísticas Implementadas}

Implementámos quatro heurísticas para utilização com os algoritmos informados:

\textbf{Euclidiana}: $h = \sqrt{(x_2-x_1)^2 + (y_2-y_1)^2}$ --- distância em linha reta, admissível em redes rodoviárias.

\textbf{Manhattan}: $h = |x_2-x_1| + |y_2-y_1|$ --- soma das diferenças absolutas, útil em grelhas urbanas regulares.

\textbf{Tempo Estimado}: $h = d_{euclidiana}/v_{media}$ --- fornece estimativa temporal relevante quando o objetivo é minimizar tempo.

\textbf{Composta}: $h = (d_{euclidiana}/v_{media}) \times f_{transito}$ --- ajusta pelo fator de trânsito atual, oferecendo maior realismo.

\section{Estratégias de Alocação}

Desenvolvemos três estratégias de alocação de veículos a pedidos, cada uma otimizando diferentes critérios:

\subsection{Veículo Mais Próximo (NearestVehicleStrategy)}

Atribui a cada pedido o veículo disponível com menor distância ao ponto de recolha. Simples e eficiente computacionalmente, minimiza tempo de resposta inicial. Não considera custos da viagem completa nem características do veículo.

\subsection{Menor Custo (LowestCostStrategy)}

Seleciona o veículo que minimiza o custo operacional total: $(d_{pickup} + d_{viagem}) \times custo_{km}$. Favorece naturalmente veículos elétricos (0.12 EUR/km) face aos de combustão (0.25 EUR/km). Adequada quando o objetivo principal é otimização financeira.

\subsection{Ecológica (EcoFriendlyStrategy)}

Prioriza veículos elétricos sempre que tenham autonomia suficiente. Primeiro filtra elétricos disponíveis com autonomia adequada e ordena por distância; se nenhum viável, considera combustão. Minimiza emissões de CO$_2$, alinhando-se com objetivos de sustentabilidade.

\section{Simulação Dinâmica}

Implementámos um simulador orientado a eventos através da classe \texttt{TaxiSimulator}. Esta arquitetura processa ocorrências por ordem temporal, sendo mais eficiente que passos temporais fixos pois apenas processa eventos quando efetivamente ocorrem.

O simulador mantém uma fila de eventos ordenada por tempo. Os tipos de eventos suportados incluem:
\begin{itemize}
    \item Chegada de novos pedidos
    \item Atribuição de veículos a pedidos
    \item Início e fim de viagens
    \item Início e fim de recargas
    \item Eventos dinâmicos (trânsito, estradas cortadas, falhas em estações)
\end{itemize}

O ciclo principal extrai o próximo evento, avança o tempo simulado, processa o evento invocando o handler apropriado e adiciona novos eventos gerados. O ciclo termina quando a fila fica vazia ou se atinge o tempo máximo.

Para simular condições realistas, implementámos mecanismos de atualização de trânsito (fatores multiplicativos nas arestas), bloqueio de estradas (arestas intransitáveis temporariamente) e falhas de estações (indisponibilidade de pontos de recarga).

\section{Resultados e Avaliação}

\subsection{Configuração Experimental}

Configurámos uma frota com 5 veículos: 3 elétricos (EV001, EV002, EV003) com autonomia de 250 km e custo de 0.12 EUR/km, e 2 a combustão (CV001, CV002) com 600 km e 0.25 EUR/km. Simulámos 30 pedidos em 180 minutos com semente aleatória fixa para reprodutibilidade.

\subsection{Métricas Quantitativas}

Os resultados nas seis métricas obrigatórias do enunciado:

\begin{center}
\begin{tabular}{|l|c|c|l|}
\hline
\textbf{Métrica} & \textbf{Valor} & \textbf{Unidade} & \textbf{Observação} \\
\hline
Tempo médio resposta & 0.14 & min & Atribuição rápida \\
Custos operacionais & 32.96 & EUR & Predominância elétricos \\
Emissões CO$_2$ & 11.92 & kg & Apenas combustão \\
Taxa atendimento & 96.67 & \% & 29/30 pedidos \\
Km vazios & 71.61 & km & 38\% do total \\
\hline
\end{tabular}
\end{center}

\subsection{Análise dos Resultados}

Os veículos elétricos apresentaram custo por pedido de 0.55-0.88 EUR versus 1.75-1.87 EUR dos de combustão --- diferença de aproximadamente 3x que confirma a vantagem económica quando a autonomia não é limitante. 100\% das emissões de CO$_2$ provêm dos veículos a combustão.

A distribuição foi assimétrica: elétricos atenderam 20 pedidos (7+5+8), combustão atenderam 9 pedidos (4+5), resultado da estratégia que prioriza menor custo. A taxa de 96.67\% indica resposta adequada à procura, com apenas 1 pedido pendente por limitações temporais.

Os 38\% de km vazios representam oportunidade de melhoria através de posicionamento proativo da frota em zonas de elevada procura.

\section{Discussão e Decisões de Design}

\subsection{Arquitetura Modular}

O sistema foi estruturado em cinco módulos independentes:
\begin{itemize}
    \item \texttt{models}: classes de domínio (Vehicle, Request, City, Station)
    \item \texttt{algorithms}: algoritmos de procura e heurísticas
    \item \texttt{strategies}: estratégias de alocação
    \item \texttt{simulation}: simulador orientado a eventos e métricas
    \item \texttt{utils}: classes auxiliares (Node, Problem)
\end{itemize}

Esta separação facilita testabilidade (cada módulo verificável isoladamente) e extensibilidade (adicionar novos algoritmos ou estratégias sem modificar código existente), seguindo o princípio Open-Closed.

\subsection{Decisões Técnicas}

Optámos por simulação orientada a eventos por ser mais eficiente que passos temporais fixos, dado que o ambiente apresenta períodos sem atividade significativa. A função de custo multi-objetivo com pesos configuráveis permite ajustar prioridades conforme objetivos operacionais.

\subsection{Tratamento de Casos Limite}

Durante o desenvolvimento, identificámos e tratámos casos limite:
\begin{itemize}
    \item Autonomia insuficiente: procura rota via estação de recarga
    \item Estações indisponíveis: pedidos entram em fila com prioridade
    \item Empate de custos: distância como critério de desempate
\end{itemize}

\section{Conclusões}

O sistema desenvolvido responde aos requisitos do enunciado: 5 algoritmos de procura (BFS, DFS, UCS, A*, Greedy), 4 heurísticas (Euclidiana, Manhattan, Tempo Estimado, Composta), 3 estratégias de alocação (NearestVehicle, LowestCost, EcoFriendly) e 6 métricas obrigatórias.

As simulações demonstraram taxa de atendimento de 96.67\% com custos de 32.96 EUR e emissões de 11.92 kg CO$_2$. O simulador orientado a eventos modelou adequadamente condições dinâmicas. Como trabalho futuro: posicionamento proativo de veículos, previsão de procura e ride-sharing.

\end{document}
